\subsubsection{开源代码逻辑、优缺点与边界}

代码流水线为：
\begin{quote}
\small\ttfamily
data -> instance -> model -> solver -> repair\\
-> evaluate -> simulate -> report
\end{quote}
Python 用 PuLP/Pyomo 或 SciPy 建立确定性基线，用
NetworkX 验证图算法接口；MATLAB 用 Optimization Toolbox/Global Optimization
Toolbox 建立第二实现。随机算法统一接收种子、预算和停止阈值，输出每代最优值、
可行率及 Pareto 集。仿真模块只通过明确的数据类接收方案，禁止在仿真循环内
偷偷改变约束。

精确规划的优势是可解释、可复算并可给上下界，缺点是非凸或大规模整数问题可能
耗时；启发式的优势是能处理混合和非光滑结构，缺点是没有全局最优保证且对参数
敏感。模型边界由结构而非题型标签决定：可线性化的约束先线性化；不确定性已知
分布可用随机规划或 Monte Carlo，只有区间信息时优先稳健情景；动态图路径与
静态最短路分开处理。

